\begin{abstract}
Kameyama--Nawata~\cite{Kameyama17}（以下 KN）は，refined Chern--Simons 不変量を resolved conifold 上の refined BPS 不変量に書き換え，
それらが非負整数になるという正値性予想を，framing 0 のトーラス結び目・トーラス絡み目と，いくつかの thin な結び目で確かめた．
本論文では，この枠組みを framing をもつ結び目に拡張し，KN が用いたデータを二つの方向に超えて，
$G$ 展開の前の係数 $\hbN$ と後の refined BPS 不変量 $\hN$ の正整数性を検証する．

まず，framing $x$ の不変量を $\rCS^{(x)}_\lambda=\fr[\lambda]^x\rCS_\lambda$ で定める．
$\fr[\lambda]$ は Taki の refined topological vertex の framing の因子（の逆数）に一致し，$t=q$ で Mariño--Vafa の framing の規則に戻る．
$\bt^{-|\rho|x}\hf^{(x)}_\rho$ は $x$ の指数和であり，KN の枠組みの構造的な性質（転置対称性，Laurent 性，$G$ 展開が閉じる条件）は $x$ によらない．
Mariño--Vafa の符号 $(-1)^{|\lambda|x}$ は，Laurent 性を保つために必要であることを示す．

次に，結び目と framing の両方の方向で検証する．
(i) ホモロジーを入力とする限り，任意の結び目・任意の framing で基本表現 $\rho=\yng(1)$ の正値性が成り立つ（命題）．11 交点以下の全素結び目 801 個で確かめた．
(ii) 2 箱では，Laurent 性の必要条件 $\dim H^{\yng(2)}=(\dim H^{\yng(1)})^2$ を導き，予想式を入力として，ツイスト結び目の全ての撚りの数での正値性を示す．
また 8 個の結び目では，全ての整数の framing で正値性を示す．
(iii) 3 箱では，トーラス結び目 7 個と代数的な iterated torus 結び目 7 個（3 本ケーブルと 2 段の反復を含む）について，
定義どおりの refined 不変量から，$|\rho|\le 3$ の全ての $\rho$ で $G$ 展開の前後とも正であることを計算で確かめた．
最大の例は約 183 万項である．
これらに基づき，framed knot の refined LMOV 不変量の正値性を予想として提出する．
\end{abstract}
